\notebook{09-the-markov-state}
The questions of this section concern the decayed counts $\decayedCounts$
of equation~\eqref{notes:eq.decayedCounts},
the Markov property of Theorem~\ref{notes:thm.markovState}
and its homogeneity in time,
and their use in the forward means of Proposition~\ref{notes:prop.forwardMean}.
The two types are the market buys and the market sells,
the first and the second component of every vector,
numbered 0 and 1 by the package.
\codehl{Clock} and \codehl{decayed\_counts\_at} are imported from \codehl{unito26.lob.hawkes};
a \codehl{Clock} is an array of times at which a state is read.
For a specification \codehl{params},
\codehl{params.forward\_mean\_state(z, h)} is $m_{s,s+h}(z)$
of equation~\eqref{notes:eq.forwardMeanState},
and \codehl{params.forward\_mean\_counts(z, h)} is the right-hand side
of equation~\eqref{notes:eq.forwardMeanCounts} at $t = s + h$,
each a vector with a component per type.
\codehl{ExponentialHawkes(params, rng)} simulates a path started empty at the time 0,
$\multiCountingProc(0) = 0$ and $\decayedCounts(0+) = 0$.
\codehl{simulator.events(T)} yields the pairs \codehl{(time, type)} of its events
up to the time $T$, in order,
and a later call continues the same path.
The one-way pair, \codehl{config.one\_way\_pair\_params()}, has
\begin{equation*}
 \baseIntensity = \begin{pmatrix} 0.4 \\ 0.1 \end{pmatrix},
 \qquad
 \excitation = \begin{pmatrix} 0.8 & 0 \\ 4 & 2.4 \end{pmatrix},
 \qquad
 \decay = 4 ,
\end{equation*}
so that a buy excites both types and a sell excites only the sells.

\begin{question}
\label{question.restartTheStateAfterAMarketSell}
The package reads the state of a path on a clock with the function below,
in which \codehl{state} starts as a copy of \codehl{initial\_state},
the state $\decayedCounts(0+)$.
\begin{lstlisting}[caption={\protect\codehl{unito26.lob.hawkes.decayed\_counts\_at}, abridged.}, label=listing.restartTheStateAfterAMarketSellQuote]
def decayed_counts_at(
    decay: Decay,
    times: np.ndarray,
    types: np.ndarray,
    initial_state: np.ndarray,
    clock: Clock,
) -> np.ndarray:
# ...
    beta = float(decay)
# ...
    last = 0.0
    event = 0
    out = np.empty((clock.size, state.size))
    for k, now in enumerate(clock):
        while event < times.size and times[event] < now:
            state *= np.exp(-beta * (times[event] - last))
            state[types[event]] += 1.0
            last = times[event]
            event += 1
        out[k] = state * np.exp(-beta * (now - last))
    return out
\end{lstlisting}
A path of four market orders is read at the times 3 and 5 from its whole history,
and again from its second order, a market sell at $s = 2$:
the events kept are shifted by $s$, as is the clock,
and the reading starts from a state \codehl{z} at $s$.
With $\decay = \log 2$, $e^{-\decay} = 1/2$.
\begin{lstlisting}[caption={Four restarts at a market sell}, label=listing.restartTheStateAfterAMarketSellRun]
beta = np.log(2.0)
times = np.array([1.0, 2.0, 3.0, 4.0])
types = np.array([0, 1, 1, 0])
clock = Clock(np.array([3.0, 5.0]))
zero = np.zeros(2)
k = 1
s = times[k]
whole = decayed_counts_at(beta, times, types, zero, clock)
before = decayed_counts_at(beta, times, types, zero, Clock([s]))[0]
after = before + np.eye(2)[types[k]]

def restart(z, keep):
    return decayed_counts_at(
        beta, times[keep] - s, types[keep], z, Clock(clock - s)
    )

for z, keep in ((before, times >= s), (before, times > s),
                (after, times > s), (after, times >= s)):
    print(np.allclose(restart(z, keep), whole))
\end{lstlisting}
What does the snippet print, line by line?
\begin{choices}
 \choice \codehl{False}, \codehl{False}, \codehl{True}, \codehl{False}
 \choice \codehl{False}, \codehl{True}, \codehl{False}, \codehl{False}
 \correctchoice \codehl{True}, \codehl{False}, \codehl{True}, \codehl{False}
 \choice \codehl{True}, \codehl{True}, \codehl{True}, \codehl{True}
\end{choices}
\explanation{The strict inequality reads $\decayedCounts$ before the jump of an event on the clock,
so \codehl{before} is $\decayedCounts(s) = (0.5, 0)$
and \codehl{after} is $\decayedCounts(s+) = (0.5, 1)$.
A restart reproduces the whole history when it counts the sell at $s$ exactly once:
in the state, or as the event at the shifted time 0,
which precedes the shifted clock and is read.}
\end{question}

\begin{question}
\label{question.millisecondClockKeepsTheState}
A feed stamps events in milliseconds after the open,
and the one-way pair is a specification in seconds.
The snippet below writes the same specification on the clock in milliseconds,
reads the state at one second after three market orders,
and forecasts the market orders of the next half second.
\begin{lstlisting}[caption={The state and the forecast on two clocks}, label=listing.millisecondClockKeepsTheState]
seconds = config.one_way_pair_params()
millis = HawkesParams(
    baseline=seconds.baseline / 1000,
    excitation=seconds.excitation / 1000,
    decay=seconds.decay / 1000,
)
times = np.array([0.2, 0.35, 0.9])
types = np.array([0, 1, 1])
zero = np.zeros(2)
z = decayed_counts_at(seconds.decay, times, types, zero, Clock([1.0]))[0]
z_ms = decayed_counts_at(
    millis.decay, 1000 * times, types, zero, Clock([1000.0])
)[0]
expected = seconds.forward_mean_counts(z, 0.5)
print(np.allclose(z_ms, z / 1000), np.allclose(z_ms, z))
print(np.allclose(millis.forward_mean_counts(z_ms, 0.5), expected),
      np.allclose(millis.forward_mean_counts(z_ms, 500.0), expected))
\end{lstlisting}
What does the snippet print?
\begin{choices}
 \correctchoice \codehl{False True}, then \codehl{False True}.
 \choice \codehl{True False}, then \codehl{False True}.
 \choice \codehl{False True}, then \codehl{True False}.
 \choice \codehl{False True}, then \codehl{False False}.
\end{choices}
\explanation{By Remark~\ref{notes:remark.timeRescaling} the flow on the new clock
has the decayed counts $\decayedCounts(t/c)$, with $c = 1000$:
each term $e^{-\decay(t - \arrivalTimes_{j})}$ is unchanged when the times are multiplied by $c$
and the decay divided by it.
The intensities are rates and are divided by $c$;
the expected counts are numbers of events, the same over the same half second,
which on the new clock is a horizon of 500.}
\end{question}

\begin{question}
\label{question.totalIntensityIsNotTheState}
A desk summarises the flow by the total intensity
$\totalIntensity = \mathbf{1}\transpose(\baseIntensity + \excitation\decayedCounts)$,
which sets the rate of the next market order,
and forecasts from it the number of market orders over the next ten seconds.
The snippet compares two states on each of two specifications.
The first is the self-exciting pair, with $\decay = 4$ and
\begin{equation*}
 \baseIntensity = \begin{pmatrix} 0.25 \\ 0.25 \end{pmatrix},
 \qquad
 \excitation = \begin{pmatrix} 2.8 & 0.2 \\ 0.2 & 2.8 \end{pmatrix};
\end{equation*}
the second is the one-way pair.
Their branching ratios are 0.75 and 0.6,
and over ten seconds $e^{-\decay(1 - \branchingRatio)h}$ is $e^{-10}$ and $e^{-16}$,
negligible at two decimals.
\begin{lstlisting}[caption={Two states of equal total intensity}, label=listing.totalIntensityIsNotTheState]
def total_intensity(params, z):
    return (params.baseline + params.excitation @ z).sum()

def total_count(params, z):
    return params.forward_mean_counts(z, 10.0).sum()

for params, z1, z2 in (
    (config.self_exciting_pair_params(), [1.0, 0.0], [0.5, 0.5]),
    (config.one_way_pair_params(), [0.5, 0.0], [0.0, 1.0]),
):
    z1, z2 = np.array(z1), np.array(z2)
    gap = total_intensity(params, z1) - total_intensity(params, z2)
    more = total_count(params, z1) - total_count(params, z2)
    print(f"{gap:.2f} {more:.2f}")
\end{lstlisting}
What does the snippet print?
\begin{choices}
 \choice \codehl{0.00 0.00} on both lines.
 \correctchoice \codehl{0.00 0.00}, then \codehl{0.00 0.19}.
 \choice \codehl{0.00 0.00}, then \codehl{0.00 -0.19}.
 \choice \codehl{0.00 0.00}, then \codehl{0.00 -0.31}.
\end{choices}
\explanation{$\totalIntensity$ weights $z$ by the column sums of $\excitation$, $(3, 3)$ and $(4.8, 2.4)$,
and the states of each pair have equal weights.
Over ten seconds the expected count weights $z$ by the expected numbers of descendants $C_{\eone} - 1$
of equation~\eqref{notes:eq.responseIntegral}, $(3, 3)$ and $(3.375, 1.5)$,
which are proportional to the column sums on the first pair only:
$0.5 \cdot 3.375 - 1.5 = 0.1875$.
The cluster sizes $C_{\eone}$, which count the ancestor too, give $-0.31$.
$\totalIntensity$ is a function of the state of Theorem~\ref{notes:thm.markovState}, not a state.}
\end{question}

\begin{question}
\label{question.fanStartsTheClockAtZero}
At $s = 30$ a desk holds the state $z = \decayedCounts(s+) = (0.5, 2)$ of the one-way pair
and estimates
$\Expectation[\multiCountingProc(s+1) - \multiCountingProc(s) \mid \decayedCounts(s+) = z]$
by the mean of independent draws of one of the four functions below.
\codehl{ExponentialHawkes.from\_state(params, z, rng)} is a simulator at the time 0 whose state is \codehl{z}.
\codehl{simulator.time} is the time the simulator has reached,
\codehl{simulator.decayed\_counts} its state,
and \codehl{events} draws each event by the method \codehl{step}, quoted below with \codehl{\_advance}.
In it \codehl{total\_intensity} is $\totalIntensity$ computed from the state,
\codehl{\_baseline\_total} is $\mathbf{1}\transpose\baseIntensity$,
and \codehl{\_draw\_type\_and\_jump} draws the type
with probabilities proportional to $\baseIntensity + \excitation\decayedCounts$,
adds its unit vector to the state, and returns it.
\begin{lstlisting}[caption={\protect\codehl{unito26.lob.hawkes.ExponentialHawkes.step}, abridged.}, label=listing.fanStartsTheClockAtZeroStep]
def step(self) -> tuple[float, int]:
    beta = self.params.decay
    excess = self.total_intensity - self._baseline_total  # D >= 0

    u1, u2 = self.rng.random(2)
# ...
    wait_baseline = -np.log(u1) / self._baseline_total
# ...
    if excess <= 0.0:
        wait_excited = np.inf
    else:
        remaining = 1.0 + beta * np.log(u2) / excess
        wait_excited = -np.log(remaining) / beta if remaining > 0.0 else np.inf

    self._advance(min(wait_baseline, wait_excited))
    return self.time, self._draw_type_and_jump()
\end{lstlisting}
\begin{lstlisting}[caption={\protect\codehl{unito26.lob.hawkes.\_HawkesState.\_advance}}, label=listing.fanStartsTheClockAtZeroAdvance]
def _advance(self, elapsed: float) -> None:
    """Decay the state forward by ``elapsed``, without an event."""
    self.time += elapsed
    self.decayed_counts *= np.exp(-self.params.decay * elapsed)
\end{lstlisting}
\begin{lstlisting}[caption={Four estimates of the expected counts after $s$}, label=listing.fanStartsTheClockAtZeroEstimates]
params = config.one_way_pair_params()
z, s, h = np.array([0.5, 2.0]), 30.0, 1.0

def counts(simulator, start):
    ahead = [e for t, e in simulator.events(start + h) if t > start]
    return np.bincount(ahead, minlength=2)

def estimate_a(rng):
    return counts(ExponentialHawkes.from_state(params, z, rng), 0.0)

def estimate_b(rng):
    return counts(ExponentialHawkes.from_state(params, z, rng), s)

def estimate_c(rng):
    return counts(ExponentialHawkes(params, rng), s)

def estimate_d(rng):
    simulator = ExponentialHawkes.from_state(params, z, rng)
    simulator.time = s
    return counts(simulator, s)
\end{lstlisting}
Which estimates tend to the expectation as the number of draws grows?
\begin{choices}
 \choice \codehl{estimate\_a} only.
 \choice \codehl{estimate\_a}, \codehl{estimate\_b} and \codehl{estimate\_d}.
 \choice \codehl{estimate\_a}, \codehl{estimate\_c} and \codehl{estimate\_d}.
 \correctchoice \codehl{estimate\_a} and \codehl{estimate\_d}.
 \choice All four.
\end{choices}
\explanation{The step reads the state and never the time, which it only accumulates,
so the path after a time has a law fixed by the state at that time:
\codehl{estimate\_d} is \codehl{estimate\_a} on a clock moved by $s$.
\codehl{estimate\_b} counts thirty seconds after the path was at $z$,
and \codehl{estimate\_c} on a path that was never set to $z$;
both estimate about the stationary count $h\stationaryIntensity = (0.5, 1.5)$.}
\end{question}

\begin{question}
\label{question.forecastBetweenEventsFollowsTheFlow}
A desk forecasts the market orders of the second that follows $t = 10$,
a time at which no event occurs.
It keeps the time \codehl{last} of the last event at or before $t$
and the state \codehl{state} just after it,
which \codehl{simulator.decayed\_counts} holds while \codehl{events} waits at that event.
\begin{lstlisting}[caption={Four forecasts on a timer}, label=listing.forecastBetweenEventsFollowsTheFlowForecasts]
params = config.one_way_pair_params()
t, h = 10.0, 1.0

def forecast_a(last, state):
    z = state * np.exp(-params.decay * (t - last))
    return params.forward_mean_counts(z, h)

def forecast_b(last, state):
    return params.forward_mean_counts(state, h)

def forecast_c(last, state):
    return (params.forward_mean_counts(state, t - last + h)
            - params.forward_mean_counts(state, t - last))

def forecast_d(last, state):
    z = state * np.exp(-params.decay * t)
    return params.forward_mean_counts(z, h)
\end{lstlisting}
\begin{lstlisting}[caption={The forecasts against the counts realised}, label=listing.forecastBetweenEventsFollowsTheFlowRun]
rng = np.random.default_rng(9)
forecasts = (forecast_a, forecast_b, forecast_c, forecast_d)
errors = {forecast: [] for forecast in forecasts}
for _ in range(3000):
    simulator = ExponentialHawkes(params, rng)
    last, state = 0.0, np.zeros(2)
    for time, _ in simulator.events(t):
        last, state = time, simulator.decayed_counts.copy()
    ahead = [event_type for _, event_type in simulator.events(t + h)]
    realised = np.bincount(ahead, minlength=2)
    for forecast in forecasts:
        errors[forecast].append(realised - forecast(last, state))
\end{lstlisting}
For which forecasts does the mean of \codehl{errors[forecast]} tend to zero in both components
as the number of paths grows?
\begin{choices}
 \correctchoice \codehl{forecast\_a} only.
 \choice \codehl{forecast\_a} and \codehl{forecast\_c}.
 \choice \codehl{forecast\_c} only.
 \choice \codehl{forecast\_a} and \codehl{forecast\_b}.
\end{choices}
\explanation{Between events the state follows the flow $e^{-\decay u}z$ of Theorem~\ref{notes:thm.markovState},
so \codehl{forecast\_a} conditions on $\decayedCounts(t+)$, the history at $t$.
\codehl{forecast\_c} conditions on the state after the last event alone,
and so averages over paths with further events before $t$,
which this path has not had: it overstates the sells.
\codehl{forecast\_b} omits the flow, and \codehl{forecast\_d} runs it from the open.}
\end{question}

\begin{question}
\label{question.forecastAfterAMarketSellUsesThePostJumpState}
A forecast of the market orders in the second that follows an event
is checked on simulated paths of the one-way pair.
The event is the first market sell of each path, at $s$,
and the forecast is made from the state read at $s$,
with or without the unit vector of the sell.
The value $e^{-1.6} \approx 0.20$ may be used.
\begin{lstlisting}[caption={The error of a forecast made at the $k$-th event}, label=listing.forecastAfterAMarketSellUsesThePostJumpStateError]
params = config.one_way_pair_params()
h = 1.0

def forecast_error(times, types, k, post_jump):
    s = times[k]
    z = decayed_counts_at(params.decay, times, types, np.zeros(2),
                          Clock(times[k:k + 1]))[0]
    if post_jump:
        z = z + np.eye(2)[types[k]]
    ahead = (times > s) & (times <= s + h)
    realised = np.bincount(types[ahead], minlength=2)
    return realised - params.forward_mean_counts(z, h)

def path(rng, horizon):
    events = list(ExponentialHawkes(params, rng).events(horizon))
    times = np.array([time for time, _ in events])
    types = np.array([event_type for _, event_type in events])
    return times, types
\end{lstlisting}
\begin{lstlisting}[caption={The errors at the first market sell}, label=listing.forecastAfterAMarketSellUsesThePostJumpStateRun]
rng = np.random.default_rng(7)
errors = {False: [], True: []}
for _ in range(2000):
    times, types = path(rng, 30.0)
    k = np.flatnonzero(types == 1)[0]
    for post_jump in errors:
        error = forecast_error(times, types, k, post_jump)
        errors[post_jump].append(error)
\end{lstlisting}
Every path has its first sell before the time 29.
Which statement holds of the mean errors, realised less forecast?
\begin{choices}
 \choice Both forecasts are unbiased: the mean errors are close to zero in both components.
 \correctchoice The forecast without the jump understates the sells by about 1.2 and not the buys;
  the forecast with it is unbiased.
 \choice The forecast without the jump understates both the buys and the sells;
  the forecast with it is unbiased.
 \choice The forecast without the jump is unbiased;
  the forecast with it overstates the sells, counting the sell at $s$ twice.
\end{choices}
\explanation{\codehl{decayed\_counts\_at} reads $\decayedCounts(s)$, before the jump,
and the forward mean is conditional on $\decayedCounts(s+) = \decayedCounts(s) + \unitVector_{2}$,
Remark~\ref{notes:remark.prePostJump}.
Without the jump the forecast misses the descendants in $(s, s+1]$ of the sell at $s$,
$\excitation\meanResponseIntegral(1)\unitVector_{2}$.
The sells excite only the sells, so $\hurwitzMatrix\unitVector_{2} = -1.6\,\unitVector_{2}$
and this is $(0, 2.4 \cdot (1 - e^{-1.6})/1.6) \approx (0, 1.2)$.}
\end{question}

\begin{question}
\label{question.forecastFromAStoppingTime}
The check of Listing~\ref{listing.forecastAfterAMarketSellUsesThePostJumpStateError}
is repeated with the state after the jump,
and with four rules for the event at which the forecast is made,
each returning its index in the path.
\begin{lstlisting}[caption={Four rules for the time of the forecast}, label=listing.forecastFromAStoppingTimeRules]
def first_after_ten(times, types):
    return np.flatnonzero(times > 10.0)[0]

def first_after_a_lull(times, types):
    lull = np.diff(times, prepend=0.0) > 0.5
    return np.flatnonzero((times > 10.0) & lull)[0]

def first_before_a_lull(times, types):
    lull = np.diff(times, append=np.inf) > 0.5
    return np.flatnonzero((times > 10.0) & lull)[0]

def largest_state(times, types):
    zero = np.zeros(2)
    z = decayed_counts_at(params.decay, times, types, zero, Clock(times))
    total = (z + np.eye(2)[types]).sum(axis=1)
    return np.argmax(total[times < 39.0])
\end{lstlisting}
\begin{lstlisting}[caption={The errors at the four times}, label=listing.forecastFromAStoppingTimeRun]
rng = np.random.default_rng(8)
rules = (first_after_ten, first_after_a_lull,
         first_before_a_lull, largest_state)
errors = {rule: [] for rule in rules}
for _ in range(2000):
    times, types = path(rng, 40.0)
    for rule in rules:
        k = rule(times, types)
        errors[rule].append(forecast_error(times, types, k, True))
\end{lstlisting}
On every path each rule returns an event before the time 39.
For which rules does the mean error tend to zero in both components
as the number of paths grows?
\begin{choices}
 \choice All four.
 \choice \codehl{first\_after\_ten} only.
 \choice \codehl{first\_after\_ten}, \codehl{first\_after\_a\_lull} and \codehl{first\_before\_a\_lull}.
 \choice \codehl{first\_after\_ten}, \codehl{first\_after\_a\_lull} and \codehl{largest\_state}.
 \correctchoice \codehl{first\_after\_ten} and \codehl{first\_after\_a\_lull}.
\end{choices}
\explanation{The law of the path after $s$ is fixed by $\decayedCounts(s+)$
at a fixed time and at a stopping time,
whose occurrence is decided by the path up to it.
The gap before an event is known when it arrives.
\codehl{first\_before\_a\_lull} picks events after which nothing arrives for half a second,
and \codehl{largest\_state} the event after which the state is never as high again before the time 39:
both look ahead, and the flow after them is quieter than the forecast.}
\end{question}

\begin{question}
\label{question.forecastAWindowThatOpensLater}
At $s = 30$ a desk holds the state $z = \decayedCounts(s+) = (0.5, 2)$ of the one-way pair
and wants the expected market buys and sells of the second second from now,
$\Expectation[\multiCountingProc(s+2) - \multiCountingProc(s+1) \mid \decayedCounts(s+) = z]$.
\begin{lstlisting}[caption={Four forecasts of a window that opens one second ahead}, label=listing.forecastAWindowThatOpensLater]
params = config.one_way_pair_params()
s, z = 30.0, np.array([0.5, 2.0])

def window_a(z):
    return (params.forward_mean_counts(z, 2.0)
            - params.forward_mean_counts(z, 1.0))

def window_b(z):
    m = params.forward_mean_state(z, 1.0)
    return params.forward_mean_counts(m, 1.0)

def window_c(z):
    return params.forward_mean_counts(z, 1.0)

def window_d(z):
    return (params.forward_mean_counts(z, s + 2.0)
            - params.forward_mean_counts(z, s + 1.0))
\end{lstlisting}
Which functions return it?
\begin{choices}
 \choice \codehl{window\_a} only.
 \correctchoice \codehl{window\_a} and \codehl{window\_b}.
 \choice \codehl{window\_a}, \codehl{window\_b} and \codehl{window\_d}.
 \choice \codehl{window\_a} and \codehl{window\_d}.
 \choice All four.
\end{choices}
\explanation{By Proposition~\ref{notes:prop.forwardMean} the expectation over $(s, s+h]$
depends on $h$ alone,
and \codehl{window\_a} differences the horizons 2 and 1.
\codehl{window\_b} conditions on the state at $s+1$ and averages over it:
the forecast is affine in the state,
so the average is the forecast from the mean state $m_{s,s+1}(z)$.
\codehl{window\_d} measures the horizons from the time 0, and gives about $\stationaryIntensity$;
\codehl{window\_c} forecasts the first second.}
\end{question}

\begin{question}
\label{question.liveStateByRecursionAgainstRecomputation}
A live forecaster needs the state $\decayedCounts(\arrivalTimes[ ]_{k}+)$ after every event of a stream.
The two functions below return the same array, a row per event.
\codehl{np.bincount(types, weights, minlength=2)} returns, for each type,
the sum of the weights of its events.
Each function is timed on the first $n$ events of one path of the one-way pair,
for $n = 1000, 2000, 4000, 8000, 16000$.
\begin{lstlisting}[caption={Two ways to the state after every event}, label=listing.liveStateByRecursionAgainstRecomputation]
def live_states_a(decay, times, types):
    state, last, out = np.zeros(2), 0.0, []
    for now, event_type in zip(times, types):
        state = state * np.exp(-decay * (now - last))
        state[event_type] += 1.0
        last = now
        out.append(state)
    return np.array(out)

def live_states_b(decay, times, types):
    out = []
    for k in range(len(times)):
        weights = np.exp(-decay * (times[k] - times[: k + 1]))
        out.append(np.bincount(types[: k + 1], weights, minlength=2))
    return np.array(out)
\end{lstlisting}
Which statement holds of the time per event?
\begin{choices}
 \choice B is faster at every $n$, its exponentials being computed by one call of numpy.
 \choice Both are about the same at every $n$, and A is faster by a constant factor.
 \choice Both grow several-fold from $n = 1000$ to $n = 16000$.
 \correctchoice That of A is about the same at every $n$,
  that of B grows several-fold from $n = 1000$ to $n = 16000$,
  and at $n = 16000$ A is faster by more than ten times.
\end{choices}
\explanation{A carries the state by the update~\eqref{notes:eq.update_of_exponential_intensity},
a fixed amount of work per event.
B evaluates~\eqref{notes:eq.decayedCounts} afresh at the $k$-th event over the $k$ events so far,
$O(n^2)$ in all;
numpy removes the interpreter's loop over the $k$ terms, not the terms.}
\end{question}
